\documentclass[12pt, ukrainian]{article}
\setlength\textwidth{170mm} \setlength\hoffset{-25mm}
%\setlength\textheight{277mm} \setlength\voffset{-38mm}
%\setlength\parindent{0mm}
\pagestyle{empty}
\usepackage[T2A]{fontenc}
\usepackage[utf8]{inputenc}
\usepackage[ukrainian]{babel}
\usepackage{graphicx}
\usepackage{caption}
\DeclareCaptionFormat{nocaption}{}
\captionsetup{format=nocaption,aboveskip=0pt,belowskip=0pt}
\usepackage[Export]{adjustbox}
\adjustboxset{max size={0.9\linewidth}{0.9\paperheight}}
\begin{document}
{20.12.2022 Варіант \No \textbf { 820 }\par}
{
\begin {Large}

{\begin {center}\textbf{Завдання 1}\end {center}}\par
Що буде виведено за виконання? \par

Якщо має бути виведено дійсне значення та інтерпретатор має здійснювати його округлення, 
то у відповіді достатньо правильно вказати два дробових розряди.

Повідомлення інтерпретатора про помилку позначати \textbf{error}.
Якщо щось буде виведено до того, як інтерпретатор видасть повідомлення про помилку, 
то воно має вказуватись у відповіді. \par

\begin{center}
	\adjustimage{max size={0.9\linewidth}{0.9\paperheight}}
	{../pict/820_1.png}
\end{center}
\newpage

{\begin {center}\textbf{Завдання 2}\end {center}}\par
Що буде виведено за виконання? \par

Якщо має бути виведено дійсне значення та інтерпретатор має здійснювати його округлення, 
то у відповіді достатньо правильно вказати два дробових розряди.

Повідомлення інтерпретатора про помилку позначати \textbf{error}.
Якщо щось буде виведено до того, як інтерпретатор видасть повідомлення про помилку, 
то воно має вказуватись у відповіді. \par

\begin{center}
	\adjustimage{max size={0.9\linewidth}{0.9\paperheight}}
	{../pict/820_2.png}
\end{center}
\newpage

{\begin {center}\textbf{Завдання 3}\end {center}}\par
Що буде виведено за виконання? \par

Якщо має бути виведено дійсне значення та інтерпретатор має здійснювати його округлення, 
то у відповіді достатньо правильно вказати два дробових розряди.

Повідомлення інтерпретатора про помилку позначати \textbf{error}.
Якщо щось буде виведено до того, як інтерпретатор видасть повідомлення про помилку, 
то воно має вказуватись у відповіді. \par

\begin{center}
	\adjustimage{max size={0.9\linewidth}{0.9\paperheight}}
	{../pict/820_3.png}
\end{center}
\newpage

{\begin {center}\textbf{Завдання 4}\end {center}}\par
Що буде виведено за виконання? \par

Якщо має бути виведено дійсне значення та інтерпретатор має здійснювати його округлення, 
то у відповіді достатньо правильно вказати два дробових розряди.

Повідомлення інтерпретатора про помилку позначати \textbf{error}.
Якщо щось буде виведено до того, як інтерпретатор видасть повідомлення про помилку, 
то воно має вказуватись у відповіді. \par

\begin{center}
	\adjustimage{max size={0.9\linewidth}{0.9\paperheight}}
	{../pict/820_4.png}
\end{center}
\newpage

{\begin {center}\textbf{Завдання 5}\end {center}}\par
Що буде виведено за виконання? \par

Якщо має бути виведено дійсне значення та інтерпретатор має здійснювати його округлення, 
то у відповіді достатньо правильно вказати два дробових розряди.

Повідомлення інтерпретатора про помилку позначати \textbf{error}.
Якщо щось буде виведено до того, як інтерпретатор видасть повідомлення про помилку, 
то воно має вказуватись у відповіді. \par

\begin{center}
	\adjustimage{max size={0.9\linewidth}{0.9\paperheight}}
	{../pict/820_5.png}
\end{center}
\newpage

{\begin {center}\textbf{Завдання 6}\end {center}}\par
Що буде виведено за виконання? \par

Якщо має бути виведено дійсне значення та інтерпретатор має здійснювати його округлення, 
то у відповіді достатньо правильно вказати два дробових розряди.

Повідомлення інтерпретатора про помилку позначати \textbf{error}.
Якщо щось буде виведено до того, як інтерпретатор видасть повідомлення про помилку, 
то воно має вказуватись у відповіді. \par

\begin{center}
	\adjustimage{max size={0.9\linewidth}{0.9\paperheight}}
	{../pict/820_6.png}
\end{center}
\newpage

{\begin {center}\textbf{Завдання 7}\end {center}}\par
Що буде виведено за виконання? \par

Якщо має бути виведено дійсне значення та інтерпретатор має здійснювати його округлення, 
то у відповіді достатньо правильно вказати два дробових розряди.

Повідомлення інтерпретатора про помилку позначати \textbf{error}.
Якщо щось буде виведено до того, як інтерпретатор видасть повідомлення про помилку, 
то воно має вказуватись у відповіді. \par

\begin{center}
	\adjustimage{max size={0.9\linewidth}{0.9\paperheight}}
	{../pict/820_7.png}
\end{center}
\newpage

{\begin {center}\textbf{Завдання 8}\end {center}}\par
Що буде виведено за виконання? \par

Якщо має бути виведено дійсне значення та інтерпретатор має здійснювати його округлення, 
то у відповіді достатньо правильно вказати два дробових розряди.

Повідомлення інтерпретатора про помилку позначати \textbf{error}.
Якщо щось буде виведено до того, як інтерпретатор видасть повідомлення про помилку, 
то воно має вказуватись у відповіді. \par

\begin{center}
	\adjustimage{max size={0.9\linewidth}{0.9\paperheight}}
	{../pict/820_8.png}
\end{center}
\newpage

{\begin {center}\textbf{Завдання 9}\end {center}}\par
Що буде виведено за виконання? \par

Якщо має бути виведено дійсне значення та інтерпретатор має здійснювати його округлення, 
то у відповіді достатньо правильно вказати два дробових розряди.

Повідомлення інтерпретатора про помилку позначати \textbf{error}.
Якщо щось буде виведено до того, як інтерпретатор видасть повідомлення про помилку, 
то воно має вказуватись у відповіді. \par

\begin{center}
	\adjustimage{max size={0.9\linewidth}{0.9\paperheight}}
	{../pict/820_9.png}
\end{center}
\newpage

{\begin {center}\textbf{Завдання 10}\end {center}}\par
Що буде виведено за виконання? \par

Якщо має бути виведено дійсне значення та інтерпретатор має здійснювати його округлення, 
то у відповіді достатньо правильно вказати два дробових розряди.

Повідомлення інтерпретатора про помилку позначати \textbf{error}.
Якщо щось буде виведено до того, як інтерпретатор видасть повідомлення про помилку, 
то воно має вказуватись у відповіді. \par

\begin{center}
	\adjustimage{max size={0.9\linewidth}{0.9\paperheight}}
	{../pict/820_10.png}
\end{center}
\newpage

\end {Large}}


\end{document}

